%% =====================================================================
%%  core/config/espaciado.tex  --  Interlineado, sangria, parrafos y notas
%% ---------------------------------------------------------------------
%%  Estilo oficio (como el machote): parrafos justificados, sangria de
%%  primera linea de 1.27 cm y aire entre parrafos. El interlineado se
%%  declara con la misma escala "tipo Word" del AWF.
%% =====================================================================

\RequirePackage{setspace}

% \setinterlineado{1|1.15|1.5|2}
\newcommand{\setinterlineado}[1]{\setstretch{#1}}

% \setsangriaparrafo{dim} / \setespacioparrafos{dim}
\newcommand{\setsangriaparrafo}[1]{\setlength{\parindent}{#1}}
\newcommand{\setespacioparrafos}[1]{\setlength{\parskip}{#1}}

% --- Defaults del oficio ---------------------------------------------
\setinterlineado{1.15}
\setsangriaparrafo{1.27cm}
%% 4pt (antes 6pt) -- Edison 2026-09-07: el aire entre párrafos
%% «ocupa demasiado». La sangría de 1.27 cm ya separa los párrafos;
%% además este \parskip se SUMA a cada \vspace de la cabecera, así
%% que bajarlo comprime todo el bloque superior. Ajustable por
%% documento con fmt_espacioparrafos.
\setespacioparrafos{4pt}

% --- Notas al pie: nunca se parten (2026-10-06) ---------------------
%% Edison 2026-10-06, en el PL-2026-009: la nota «se repartió a dos
%% páginas». LaTeX parte por defecto una nota larga entre hojas
%% (\interfootnotelinepenalty = 100) y deja su resto al pie de la
%% siguiente, lejos de la llamada. Con 10000 la nota viaja entera: si
%% no cabe, la línea con la llamada pasa a la hoja siguiente junto con
%% ella. Un documento oficial no lleva notas partidas.
\interfootnotelinepenalty=10000\relax
